\cotefile{40}

\scanpage{2}

% archive sheet number: 1

On se place dans la catégorie duale de la catégorie des algèbres $B$ (commutatives avec unité) sur l'anneau comm. $A$.

% slanted
\margin{\uline{Calculs \inserted{affines} sur} des fibrés \uline{principaux} \uline{affines}}

Une telle $B$ est un groupe de la catégorie si \unsure{c'est|l'on} une \uline{$A$-algèbre commutative unitaire} s'est donné un \uline{hom. d'algèbres}

\[
\unsure{p|\rho}_{B}\colon B\longrightarrow B\otimes_A B
\]

(``application diagonale'') qui soit ``\uline{associative}'', pour laquelle il \uline{existe} une \uline{augmentation} \inserted{\uline{d'algèbre}} $\varepsilon:B\to A$ \deleted{telle que} telle que

\[
\delta_B x=p_Bx-\deleted{\unsure{( |x}}(x\otimes1+1\otimes x)\in\unsure{J|I}\otimes J
\]

(où $J=\operatorname{Ker}\varepsilon$) --- cette augmentation est alors unique --- et un \uline{hom d'algèbres} $x\longrightarrow\check{x}$ satisfaisant à la condition

\[
x\in B,\qquad p_Bx=\sum y_i z_i\quad\Longrightarrow\quad \sum\check{y}_i z_i=\sum y_i\check{z}_i=\varepsilon(x).1
\]

--- cet hom d'algèbres est alors unique, et est une involution. De plus on aura

\[
\varepsilon(\check{x})=\varepsilon(x)
\]

\[
\text{\deleted{\unsure{et}} si}\qquad p_B\unsure{'}x=\sum y_i\otimes z_i,\quad\text{alors}\quad p_B\check{x}=\sum_i\check{z}_i\otimes\check{y}_i
\]

\uline{Si} $B$ est libre de rang fini sur $B$, alors\deleted{\unsure{s}} ces structures se reflètent directement dans $U$ par la donnée \deleted{d'une \uline{loi d'algèbre} (\unsure{non|i.e.} nécessairement commutative) \uline{associative unitaire} sur $A$, \inserted{dont la duale est une} telle qu'il existe un \unsure{hom} \unsure{$x\to\check{x}$|$x\mapsto\check{x}$}}
\scanpage{3}

% archive sheet number: 2

D'une loi de \uline{coalgèbre} \inserted{\uline{unitaire}} sur $A$, \deleted{\unsure{de|sur} $U$,} donnée par

\[
\begin{aligned}
\Delta_U &: U\unsure{\longrightarrow|-}U\times_A U\\
\unsure{J|I|\eta}_U &: U\longrightarrow A\qquad \text{augmentation (\unsure{surjective|surjection})}
\end{aligned}
\]

\deleted{compatibles en le sens qu'on \unsure{aura} si \unsure{une|d'une} augmentation}

\[
\deleted{\delta_U x=\Delta_U x-1}
\]

qui soit \uline{commutative}, \uline{associative}, d'une \uline{structure de $A$-algèbre} sur $U$ qui soit \uline{associative}, \uline{unitaire} (pas nécessairement commutative) et telle que $U\otimes_A U\to U$ soit un hom de coalgèbres unitaires [ou, ce qui revient au \unsure{même}, $\Delta_U$ est un hom d'algèbres associatives unitaires]. Enfin on suppose qu'il existe un hom \inserted{de \uline{coalgèbre} \uline{unitaire}} $x\to\check{x}$, ayant la propriété suivante :

\[
x\in U,\quad \Delta_U x\unsure{=|>}\sum y_i\otimes z_{\deleted{\unsure{j}}i}
\quad\Longrightarrow\quad \sum\deleted{\unsure{z_i\check{y}_i}}y_i\check{z}_i=\sum\deleted{\unsure{\check{z}_i y_i}}\check{y}_i z_i=J_U(x).1_U.
\]

Alors $x\to\check{x}$ est un antiautomorphisme involutif de la structure d'algèbre de $U$, et \deleted{\unsure{cela}} il est défini de façon unique par les conditions énoncées.

Soit maintenant $C$ une $A$-algèbre. On veut expliquer \unsure{que|ce que} signifie que $C$ admet $B$ comme ``groupe d'opérateurs à \deleted{droite} gauche'' dans la catégorie $\unsure{\mathcal{C}|\mathscr{C}}$. Cela signifie ceci :
\scanpage{4}

% archive sheet number: 3

On se donne un \uline{hom. d'algèbres unitaires}

\[
q_C:C\xrightarrow{\deleted{\unsure{a}}}B\otimes_A C
\]

satisfaisant aux conditions suivantes :

\uline{Associativité} On a commutativité \unsure{de|du}

\diagram{\begin{tikzcd}[ampersand replacement=\&]
B\otimes_{\deleted{\unsure{A}}}B\otimes_{\deleted{\unsure{B|A}}}\deleted{\unsure{B}}C
 \& B\otimes_{\deleted{\unsure{A}}}C \arrow[l,"{\mathrm{id}_B\otimes q_C}"']\\
B\otimes C \arrow[u,"{\unsure{p|\rho}_B\otimes\mathrm{id}_C}"]
 \& C \arrow[l,"q_C"'] \arrow[u,"q_C"']
\end{tikzcd}}

\uline{Unité} On a commutativité \unsure{de|du}

\diagram{\begin{tikzcd}[ampersand replacement=\&]
 A\otimes C=C \& B\otimes C \arrow[l,"{\varepsilon\otimes\mathrm{id}_C}"']\\
 \& C\arrow[u,"q_C"']\arrow[ul,"\sim"]
\end{tikzcd}}

i.e. on a $q_C(y)\in1\otimes y+\unsure{J|I}\otimes C$ (où $J=\operatorname{Ker}\varepsilon$)

Supposons encore que $B$ soit libre de rang fini, de sorte qu'on peut lui substituer $U$.

Alors la donnée d'un $A$-hom $q_C:C\to B\otimes_A C'$ équivaut (quels que soient les $A$-modules $C,C'$) à un $A$-hom $q_{\unsure{U|U'|\mathcal{U}}}\deleted{\unsure{B\to C}}:U\to\operatorname{Hom}_A(C,C')$ [car $B\otimes_A C'\simeq\operatorname{Hom}_A(U,C')$]. Si $C,C'$ sont des $A$-algèbres unitaires, alors dire que $q_C$ est un hom de $A$-algèbres unitaires signifie que l'on a les conditions
\scanpage{5}

% archive sheet number: 4

\[
\left\{\begin{aligned}
&q_U(y).1_C=\unsure{J|I|\eta}_U(y).1_{\deleted{\unsure{C}}C'}\\
&x,y\in C,\ u\in U,\ \Delta_U u=\sum v_i\otimes w_i\quad\Longrightarrow\quad
\unsure{p|\rho}_U(u)(xy)=\sum(p_U(\deleted{\unsure{u}}v_i)x)(p_U(w_i)y)
\end{aligned}\right.
\]

Si enfin $C=C'$, alors les \deleted{\unsure{en ce}} axiomes d'associativité et d'unité signifient que \deleted{$U$ est un} $U\to\operatorname{Hom}_A(C,C)$ est un \uline{hom de $A$-algèbres unitaires}. Considérant \uline{alors $C$ comme un $U$-module} \inserted{\uline{unitaire}} \uline{à gauche}, les deux relations précédentes s'écrivent aussi

\[
\left\{\begin{aligned}
&u.1_C=J(u).1_C\\
&x,y\in C,\ u\in U,\ \Delta_U u=\sum v_i\otimes w_i\ \Rightarrow\
u(xy)=\sum(\unsure{u_i|v_i}x)(w_i y)
\end{aligned}\right.
\]

Nous voulons maintenant exprimer que $C$, considéré comme ayant son groupe d'opérateurs $B$ dans la catégorie $\unsure{\mathcal{C}|\mathscr{C}}$ des $A$-algèbres comm. avec unité, est \deleted{\uline{principal} de base \unsure{homogène}} \inserted{fibré} \uline{principal} \deleted{de base $D$} \inserted{\uline{homogène}} \deleted{($D\unsure{\to|=|\text{est}}$ une $A$-algèbre commutative). Cela signifie donc qu'il existe un \uline{hom de $A$-algèbres unitaires}}

\diagram{\begin{tikzcd}[ampersand replacement=\&]
\deleted{C\otimes_A C}\& \deleted{B}\arrow[l,draw=none,"{\deleted{\longleftarrow}}" description]
\end{tikzcd}}
\scanpage{6}

% archive sheet number: 5

On suppose donc \deleted{donné un hom de \uline{$A$-algèbres} \unsure{unitaires}} que \deleted{$f:D\longrightarrow C$ et un \uline{hom de $A$-algèbres unitaires}}

\[
\deleted{g_C:B\longrightarrow C\otimes_{\deleted{\unsure{A|D}}}C}
\]

\deleted{rendant commutatif le diagramme suivant}

\diagram{\begin{tikzcd}[ampersand replacement=\&]
{} \& \deleted{B\otimes C}\arrow[l] \& \deleted{C}\arrow[l,"{\deleted{q_C}}"']
\end{tikzcd}}

\deleted{ayant la propriété suivante :} l'hom composé

\diagram{\begin{tikzcd}[ampersand replacement=\&]
B\otimes C \& B\otimes C\otimes C\arrow[l,"{\mathrm{id}_B\otimes\Delta_C}"'] \& C\otimes C\arrow[l,"{q_C\otimes\mathrm{id}_C}"']
\end{tikzcd}}

[i.e. l'hom $C\otimes C\xrightarrow{(q_C,p_2)}B\otimes\deleted{\unsure{A}}C$] est un \uline{isomorphisme} (de $A$-algèbres). Il revient au \unsure{même} de dire qu'il \uline{existe un hom de $A$-algèbres unitaires} $\boxed{g_C:B\longrightarrow C\otimes C}$

tel que l'on ait commutativité des

\diagram{\begin{tikzcd}[ampersand replacement=\&]
(*)\quad C\otimes C \& \deleted{\unsure{B\otimes C\otimes C}}\arrow[l,"{(g_C,\deleted{\unsure{1}}p_2\deleted{\unsure{C}})}"'] \& B\otimes C\arrow[l,"{\deleted{\unsure{g\otimes p_2}}}"',"{\deleted{\unsure{p_2}}}"] \& C\arrow[l,"q_C"']\arrow[lll,bend left=17,"p_1"]
\end{tikzcd}}

\diagram{\begin{tikzcd}[ampersand replacement=\&]
(**)\quad B\otimes C\& C\otimes C\arrow[l,"{(q_C,p_2)}"']\& B\arrow[l,"g_C"']\arrow[ll,bend left=24,"p_1"]
\end{tikzcd}}

Nous allons \deleted{\unsure{voir}} \inserted{examiner} si sous ces conditions, il est bien exact que $C$ soit fibré principal \uuline{hom} sous $G$, de \deleted{\unsure{g}} base l'objet unité $A$ de $\unsure{\mathcal{C}|\mathscr{C}}$.

Supposons donc qu'on a un hom de $A$-algèbres \inserted{commutatives unitaires}

\[
h:D\longrightarrow C
\]

rendant commutatif les diagrammes
\scanpage{7}

% archive sheet number: 6

\diagram{\begin{tikzcd}[ampersand replacement=\&]
B\otimes C\& C\arrow[l,"q_C"']\\
C\arrow[u,"{\mathrm{pr}_2}"]\&D\arrow[l,"h"']\arrow[u,"h"']
\end{tikzcd}}

\deleted{i.e. \unsure{$B$}}

est-il vrai qu'on peut factoriser $h$ \deleted{\unsure{p}} à l'aide de $h':D\to A$ (une augmentation) \deleted{que} par la formule

\[
h(x)=h'(x).1_C\qquad ??
\]

La commutativité du diagramme ci-dessus signifie simplement que $h$ prend ses valeurs dans le sous-ensemble $C'$ de $C$ \deleted{formé des $x\in$} défini par

\[
x\in C'\quad\Longleftrightarrow\quad\deleted{\unsure{q_C x|p_2}\;1\otimes x=}\ q_C x=1_B\otimes x
\]

i.e. les \uline{éléments} \inserted{\deleted{\unsure{quasi}} \uline{relativement}} \uline{primitifs} de $C$ (\unsure{v.z|viz|p.r.|rel.} : $B$) C'est là une sous-$A$-algèbre de $C$, et on voit qu'en tout état de cause (\unsure{même} si $C$ n'est pas fibré principal)

\uline{le quotient de} $C$ \uline{par le groupe d'opérateurs} $B$ \uline{dans la catégorie} $\unsure{\mathcal{C}|\mathscr{C}}$ \uline{existe}, \uline{et est donné par la sous-algèbre} $C'$ de $C$ \uline{formé des éléments} \deleted{\unsure{quasi}} \inserted{\uline{relativement}} \uline{primitifs} de $C$ (\unsure{opér.|p.r.|par.} : $B$).
\scanpage{8}

% archive sheet number: 7

\margin{\unsure{2|$Z$|$\big\}$}}

Mais attention, ce résultat ne doit pas faire illusion : Si on se place dans la catégorie des \uline{Schémas} sur $\operatorname{Spec}(A)$, il n'implique pas l'existence du quotient dans cette catégorie, \deleted{\unsure{ni}} et si le quotient y existe, il peut être distinct de celui qu'on vient de définir.

Supposons p.ex $A={}$corps alg. clos $k$, $\deleted{\unsure{E}}C={}$algèbre \inserted{affine} d'un groupe linéaire connexe non \unsure{résoluble} $\deleted{\unsure{G}}E$, $B={}$algèbre d'un sous-groupe de Borel $\deleted{\unsure{B}}G$. Alors $E/G$ dans la catégorie des schémas existe bien et est une variété complète $\ne e$. Il en résulte que le quotient dans la \uline{catégorie des variétés affines} est réduit au point $e$.

On \uline{doit donc vérifier} \uline{qu'en l'occurrence}, \deleted{\unsure{les}} \makebox[0pt][r]{\smash{\raisebox{-.9\baselineskip}{\(\left\Vert\vphantom{\rule{0pt}{2.2\baselineskip}}\right.\)}}\hspace{.6em}}[\uline{$C$ étant principal homogène}] \deleted{les éléments} \uline{$C'$ est isomorphe \unsure{:|à} $A$} \uline{[par $\lambda\to\lambda1_C$]}.

En effet, si $x$ est semi-primitif, alors le diagramme commutatif $(*)$ donne

\[
x\otimes1_C=1_C\otimes x\qquad\text{dans}\quad C\otimes C
\]
\scanpage{9}

% archive sheet number: 8

\uline{Supposons que} $\lambda\to\lambda1_C$ soit un isomorphisme de $A$ sur un facteur direct du $A$-module $C$ [pour ceci, il suffit, comme on voit facilement, que $C$ soit un $A$-module \deleted{libre de type} \inserted{ayant} une base contenant $1_C$, p.ex. que ce soit un $A$-module libre de type fini]. Alors la condition précédente implique facilement $x=\lambda1_C$, où $\lambda\in A$ est déterminé de façon unique. Dans ce cas, on a donc bien ce qu'on voulait.

\uline{N.B.} En général, \deleted{\unsure{on}} il ne sera pas vrai que le quotient de $C$ par $B$ (ds $\unsure{\mathcal{C}|\mathscr{C}}$) soit $A$, sous les conditions indiquées. On peut prendre par exemple $C=B/\unsure{J|I}B$, où $J$ est un idéal non nul de $A$, \deleted{\unsure{Alors}} \deleted{de définir $q_C$} en prenant la structure d'algèbre quotient, et définissant $q_C:C\to B\otimes C$ i.e.

\[
q_C:B/JB\longrightarrow (B\otimes B)/J.(B\otimes B)
\]

en réduisant $\unsure{p|\rho}_B$ mod $J$. Les axiomes restent vérifiés, cependant l'algèbre \unsure{semi-primitive|des primitifs} de $C$ est $A/J$, et non $A$\ldots\ldots
\scanpage{10}

% archive sheet number: 9

Supposons maintenant que $C$ soit lui aussi un $A$-module libre (de sorte que la condition relative aux quasi-primitifs sera une conséquence de l'autre condition). Je dis que \uline{pour que} \deleted{$\varphi$}

\[
\varphi:\deleted{\unsure{B}}C\otimes_A C\longrightarrow\operatorname{Hom}_A(U,C)\tag*{\unsure{$(*)$|$(\times)$}}
\]

\uline{soit un isomorphisme}, \uline{il faut et il suffit que l'hom analogue défini par} $\varphi$

\[
\psi:C\otimes U\longrightarrow\operatorname{Hom}(C,C)\tag*{$(**)$}
\]

\uline{par la condition}

\[
\boxed{\psi(x\otimes u).y=x\,u(y)}
\]

\margin{?}

\uline{soit un isomorphisme}. [Donc, si $B$ et $C$ sont tous deux libres de type fini, on a deux caractérisations \deleted{\unsure{des fibrés principaux de groupe}} \inserted{équivalentes} \inserted{du fait que $C$ soit} \deleted{fibré} un fibré principal homogène sous $C$]

Supposons en effet que $\varphi$ soit un isomorphisme, i.e. \deleted{étendons les scalaires à $A'=C$, alors}

\[
\deleted{C^{A'}\otimes_{A'}C^{A'}\longrightarrow\operatorname{Hom}_{A'}(U^{A'},C^{A'})}
\]

\deleted{est encore un \unsure{$B$|$C$|$B^{A'}$} est ``fibré hom. principal sous $B=U'$'' dans la catégorie des $A$-algèbres comm. Faisant le produit \deleted{fibré} par $A'=C$, on trouve que \inserted{ds la catégorie des $A'$-algèbres} le fibré $C^{A'}=C\otimes_A A'$ sur $B^{A'}=U'\otimes_A A'=(U^{A'})'$ est isomorphe à $B^{A'}$ lui-même, [car on a pris $A'=C$].}
\scanpage{11}

% archive sheet number: 10

Supposons \deleted{\unsure{de|à}} \inserted{à} nouveau que $B$ soit libre de type fini sur $A$, \deleted{Alors} et exprimons le fait que \unsure{de} l'on a un fibré \uline{principal} \uuline{homogène} à l'aide de $U$, $q_U$. On doit donc exprimer d'abord que l'hom d'algèbres

\[
\varphi:C\otimes C\longrightarrow\operatorname{Hom}_A(U,C)
\]

[où la structure d'algèbre du \unsure{dernier|deuxième} mb est définie par

\diagram{\begin{tikzcd}[ampersand replacement=\&]
U\arrow[r,"{\unsure{\Delta_U|\check{\Delta}_U}}"]\arrow[rrr,bend right=15,"{v*w}"'] \& U\otimes U\arrow[r,"{v\otimes w}"]\& C\otimes C\arrow[r,"{\Delta_C}"]\& C\quad{]}
\end{tikzcd}}

qui est défini par

\[
\begin{aligned}
&\deleted{\varphi(1_C\otimes x).(u)=\unsure{J|I|\eta}_U(u)x}\\
&\deleted{\varphi(x\otimes1_C)(u)=u\,x}
\end{aligned}
\]

\[
\boxed{\varphi(x\otimes y).u=u(x)\unsure{,|.}y}
\]

\makebox[0pt][r]{\smash{\raisebox{-2.3\baselineskip}{\(\left.\vphantom{\rule{0pt}{5\baselineskip}}\right]\)}}\hspace{.65em}}est un \uuline{isomorphisme} de $C\otimes C$ sur $\operatorname{Hom}_A(U,C)$

De plus, on veut que l'hom $\lambda\to\lambda.1_C$ soit un \uline{isomorphisme} de $A$ sur l'ens des éléments semi-primitifs de $C$, i.e. tels que

\[
\underline{u.x=J(u).x}
\]
\scanpage{12}

% archive sheet number: 11

Catégorie $\unsure{\mathcal{C}|\mathscr{C}}$ avec unité à droite $e$, et produits

$G$ groupe dans $\mathcal{C}$ \qquad $\unsure{p|\rho}_G:G\times G\to G$

$E$ $\mathcal{C}$-espace : $\mathcal{C}$-groupe $G$ opérant \unsure{à|:} gauche \qquad $q_E:G\times E\to E$

On dit que $E$ est formellement hom. principal. Si l'hom

\[
G\times E\xrightarrow{(\deleted{\unsure{g}}q_E,p_2)}E\times E
\]

\margin{$x=g(x,y)y$}

\margin{\unsure{\S|$S$}}

est un \uline{isomorphisme}. Cela signifie aussi que l'on peut trouver un hom

\[
E\times E\xrightarrow{g_E}G\quad\unsure{\prime\prime|,}
\]

tel que

\diagram{\begin{tikzcd}[ampersand replacement=\&]
\text{(i)}\quad E\times E\arrow[r,"{(g,p_2)}"]\arrow[rr,bend right=20,"p_1"'] \& G\times E\arrow[r,"q_E"]\& E\quad\text{est commutatif}
\end{tikzcd}}

\diagram{\begin{tikzcd}[ampersand replacement=\&]
\text{(ii)}\quad G\times E\arrow[r,"{(q_E,p_2)}"]\arrow[rr,bend right=20,"{p_{\deleted{\unsure{2}}1}}"'] \& E\times E\arrow[r,"g_E"]\& G\quad\text{est commutatif}
\end{tikzcd}}

Si $\mathcal{C}$ est la catégorie des $A$-algèbres (commutatives avec unité), alors $G,E$ correspondent à des $A$-algèbres \deleted{\unsure{comm.}} $B,C$, et on a

\[
\left\{\begin{aligned}
p_{\deleted{\unsure{G}}B}&:B\longrightarrow B\otimes B\\
q_C&:C\longrightarrow B\otimes C
\end{aligned}\right.\tag*{(*)}
\]

et le fait qu'on ait \uline{formellement} un fibré hom. principal signifie que

\[
\varphi:C\otimes C\xrightarrow{(q_C,p_2)}B\otimes C
\]

est un isomorphisme, i.e. qu'il existe

\[
g_C:B\longrightarrow C\otimes C
\]
\scanpage{13}

% archive sheet number: 12

rendant commutatifs les diagrammes

\diagram{\begin{tikzcd}[ampersand replacement=\&]
\text{(i)}\quad C\arrow[r,"q_C"]\arrow[rr,bend right=14,"p_1"']\& B\otimes C\arrow[r,"{(g_C,p_2)}"]\& C\otimes C
\end{tikzcd}}

\deleted{\unsure{(ii)}}

\diagram{\begin{tikzcd}[ampersand replacement=\&]
\text{(ii)}\quad B\arrow[r,"g_C"]\arrow[rr,bend right=20,"p_1"']\& C\otimes C\arrow[r,"{\varphi=(q_C,p_2)}"]\& B\otimes C
\end{tikzcd}}

D'autre part, les données \unsure{$(*)$|$*)$} définissent

\[
\psi:\operatorname{Hom}_A(B,C)\longrightarrow\operatorname{Hom}_A(C,C)
\]

en posant $\psi(v)=\Delta_C(v\otimes\mathrm{id}_C)q_C$ :

\diagram{\begin{tikzcd}[ampersand replacement=\&]
C\arrow[r,"q_C"]\arrow[rrr,bend right=12,"{\psi(v)}"']\& B\otimes C\arrow[r,"{v\otimes\mathrm{id}_C}"]\& \deleted{\unsure{B}}C\otimes C\arrow[r,"{\Delta_C}"]\&C
\end{tikzcd}}

Je dis que \uline{si} $\varphi$ \uline{est un isomorphisme}, $\psi$ \uline{en est un}.

On va en effet construire, utilisant $g_C$, un hom

\[
\psi':\operatorname{Hom}(C,C)\longrightarrow\operatorname{Hom}(B,C)
\]

qui sera inverse de $\psi$. On pose

\[
\psi'(w)=\Delta_C(w\otimes\mathrm{id}_C)g_C\quad :
\]

\diagram{\begin{tikzcd}[ampersand replacement=\&]
B\arrow[r,"g_C"]\arrow[rrr,bend right=13,"{\psi'(w)}"']\& C\otimes C\arrow[r,"{w\otimes\mathrm{id}_C}"]\& C\otimes C\arrow[r,"{\Delta_C}"]\& C
\end{tikzcd}}

On doit prouver $\psi'(\psi(\unsure{w|v}))=v$, $\psi\psi'(w)=w$. C'est facile :

\[
\deleted{\psi\psi'(w)=\Delta_{\unsure{C}}[(\Delta_C(w\otimes\mathrm{id}_{\unsure{C}})g_C)\otimes\mathrm{id}_C]q_C}
\]
\scanpage{14}

% archive sheet number: 13

\deleted{\unsure{$\psi'$}}

\uline{1}$^\circ$) \uline{$\psi\psi'(w)=w$}

\diagram{\resizebox{\linewidth}{!}{\begin{tikzcd}[ampersand replacement=\&,row sep=2.0em,column sep=.6em]
C\arrow[d,"{q_C}"] \&  \& C\arrow[ddd] \&  \& \deleted{\unsure{C}}C\arrow[dd] \&  \& \deleted{\unsure{C}}C\arrow[dd] \&  \& \deleted{\unsure{C|B}}C\arrow[d] \&  \& \deleted{\unsure{G|C}}C\arrow[ddd,"p_1"]\\
B\otimes C\arrow[d,shift right=2.7ex,"{g_C}"']\arrow[d,shift left=2.7ex,"{\mathrm{id}_C}"] \&  \&  \&  \&  \&  \&  \&  \& B\otimes C\arrow[dd,"{(g_{\deleted{\unsure{B}}C},p_2)}"] \& \left.\vphantom{\rule[-4.5em]{0pt}{9em}}\right\} \& \\
\rlap{$\overbrace{\phantom{C\otimes C}}$}C\otimes\underbrace{C\otimes C}\arrow[d,shift right=2.7ex,"{w}"']\arrow[d,shift left=2.7ex,"{\mathrm{id}_{C\otimes C}}"] \&  \&  \&  \& C\otimes\underbrace{C\otimes C}\arrow[dd,shift right=2.7ex,"w"']\arrow[dd,shift left=2.7ex,"{\Delta_C}"] \&  \& C\otimes\underbrace{C\otimes C}\arrow[d,shift right=2.7ex,"{\mathrm{id}_C}"']\arrow[d,shift left=2.7ex,"{\Delta_C}"] \& \left.\vphantom{\rule[-3.5em]{0pt}{7em}}\right\} \&  \&  \& \\
\underbrace{C\otimes\rlap{$\overbrace{\phantom{C\otimes C}}$}C}\otimes C\arrow[d,shift right=2.7ex,"{\Delta_C}"']\arrow[d,shift left=2.7ex,"{\mathrm{id}_C}"] \&  \& \rlap{$\deleted{\underbrace{\phantom{C\otimes C}}}$}C\otimes\underbrace{C\otimes C}\arrow[d,shift right=2.7ex,"{\mathrm{id}_C}"']\arrow[d,shift left=2.7ex,"{\Delta_C}"] \& \left.\vphantom{\rule[-3.5em]{0pt}{7em}}\right\} \&  \& \left.\vphantom{\rule[-4em]{0pt}{8em}}\right\} \& C\otimes\deleted{\unsure{B}}C\arrow[d,shift right=2.7ex,"{w}"']\arrow[d,shift left=2.7ex,"{\mathrm{id}_C}"] \&  \& C\otimes C\arrow[dd] \&  \& C\otimes C\arrow[d,"{w\otimes\mathrm{id}_C}"]\\
C\otimes C\arrow[d] \& \left.\vphantom{\rule[-3em]{0pt}{6em}}\right\} \& C\otimes C\arrow[d,"{\Delta_C}"] \&  \& C\otimes C\arrow[d] \& \deleted{\unsure{\}}} \& C\otimes C\arrow[d,"{\deleted{\unsure{\Delta_C}}}"] \&  \&  \&  \& C\otimes C\arrow[d,"{\Delta_C}"]\\
C \&  \& C \&  \& C \&  \& \deleted{C} \&  \& C \&  \& C
\end{tikzcd}}}

Sur la dernière colonne, on reconnaît aussitôt qu'on a un composé identique à $w$

2$^\circ$) $\psi'\psi\unsure{w)| (v)}=v$

\diagram{\resizebox{\linewidth}{!}{\begin{tikzcd}[ampersand replacement=\&,row sep=2.0em,column sep=.6em]
B\arrow[d,"{g_C}"] \&  \&  \&  \&  \&  \&  \&  \&  \&  \& B\arrow[dd,"p_1"]\\
C\otimes C\arrow[d,shift right=2.7ex,"{q_C}"']\arrow[d,shift left=2.7ex,"{\mathrm{id}_C}"] \&  \&  \&  \&  \&  \&  \&  \& C\otimes C\arrow[d,"{(q_C,p_2)}"] \& \left.\vphantom{\rule[-4em]{0pt}{8em}}\right\} \& \\
B\otimes\underbrace{C\otimes C}\arrow[d,shift right=2.7ex,"{v}"']\arrow[d,shift left=2.7ex,"{\mathrm{id}_{C\otimes C}}"] \&  \&  \&  \& B\otimes\underbrace{C\otimes C}\arrow[dd,shift right=2.7ex,"v"']\arrow[dd,shift left=2.7ex,"{\Delta_C}"] \&  \& B\otimes\underbrace{C\otimes C}\arrow[d,shift right=2.7ex,"{\mathrm{id}_B}"']\arrow[d,shift left=2.7ex,"{\Delta_C}"] \& \left.\vphantom{\rule[-3em]{0pt}{6em}}\right\} \& B\otimes C \&  \& B\otimes C\arrow[d,"{v\otimes\mathrm{id}_C}"]\\
\underbrace{C\otimes\rlap{$\overbrace{\phantom{C\otimes C}}$}C}\otimes C\arrow[d,shift right=2.7ex,"{\Delta_C}"']\arrow[d,shift left=2.7ex,"{\mathrm{id}_C}"] \&  \& C\otimes\underbrace{C\otimes C}\arrow[d,shift right=2.7ex,"{\mathrm{id}_C}"']\arrow[d,shift left=2.7ex,"{\Delta_C}"] \& \left.\vphantom{\rule[-3.5em]{0pt}{7em}}\right\} \&  \& \left.\vphantom{\rule[-3.5em]{0pt}{7em}}\right\} \& B\otimes C\arrow[d,shift right=2.7ex,"{v}"']\arrow[d,shift left=2.7ex,"{\mathrm{id}_C}"] \&  \&  \&  \& \deleted{\unsure{B}}C\otimes C\arrow[d,"{\Delta_C}"]\\
C\otimes C\arrow[d,"{\Delta_C}"] \& \left.\vphantom{\rule[-3.5em]{0pt}{7em}}\right\} \& \deleted{\unsure{B}}C\otimes C\arrow[d,"{\Delta_C}"] \&  \& B\otimes C \&  \& B\otimes C \&  \&  \&  \& C\\
C \&  \& C \&  \&  \&  \&  \&  \&  \&  \& 
\end{tikzcd}}}

Sur la dernière colonne, on voit aussitôt que le composé est identique à $v$
\scanpage{15}

% archive sheet number: 14

Réciproquement, supposons $\psi$ \uline{un isomorphisme}, \uline{montrons qu'alors} $\varphi$ \uline{est un isomorphisme}, du moins si on admet que $C$ admet une base finie sur $A$. Soit en effet $\psi'$ l'isomorphisme inverse de $\psi$

\[
\psi':\operatorname{Hom}_{\unsure{A|B}}(C,C)\longrightarrow\operatorname{Hom}(B,C)
\]

il s'identifie à un hom

\[
\deleted{B\otimes}\operatorname{Hom}(C,C)\otimes B\longrightarrow C
\]

i.e. à un hom

\[
\psi_1^{\deleted{\unsure{1|*}}\prime}:B\longrightarrow\operatorname{Hom}(\operatorname{Hom}(C,C),C)
\]

On veut en tirer un hom $B\to C\otimes C$ en utilisant un hom naturel

\[
\unsure{\gamma|f|\rho}:\operatorname{Hom}(\operatorname{Hom}(C,C),C)\longrightarrow C\otimes C
\]

Pour construire un tel, on note que le premier mb est isomorphe à

\[
(C'\otimes C)'\otimes C=C\otimes C'\otimes C\deleted{=C\otimes\operatorname{Hom}(C,C)}
\]

\deleted{[en général, on a \deleted{$\operatorname{Hom}($} $P\otimes\operatorname{Hom}(Q,R)\to\operatorname{Hom}(\operatorname{Hom}(P,Q),R)$ i.e. $P\otimes\operatorname{Hom}(Q,R)\otimes\operatorname{Hom}(P,Q)\to R$, donné par $(p,v,u)\to(v\circ u)(p)$]},

or l'élément unité de $C^{\unsure{*|\times}}$ définit un $A$-hom naturel

\[
\varepsilon:C'\longrightarrow A\unsure{!|,}
\]

d'où un hom $\gamma=\mathrm{id}_C\otimes\varepsilon\otimes\mathrm{id}_C$.
